\section{Experimental Setup}
\subsection{Benchmarks}
We evaluate BPGS across benchmark and stress settings that serve different roles in the paper. First, we use the standard NYUv2 multi-task benchmark \cite{silberman2012nyuv2} for the main dense-prediction comparison. The benchmark contains 795 training images and 654 validation images, with three tasks: semantic segmentation, depth estimation, and surface-normal prediction. The main NYUv2 comparison uses a 120-epoch budget and reports final-epoch metrics aggregated over seeds.

Second, we run a separate NYUv2 ablation study on a fixed 50\% training subset while keeping the validation split unchanged. This study isolates the role of the bounded chart and the initialization rule by comparing four BPGS variants against Kendall uncertainty weighting: the canonical batch-aware auto-calibrated form, a batch-aware fixed-initialization form, a stateless fixed-initialization form, and a stateless auto-calibrated form. The ablation runs use a 60-epoch budget and the same three training seeds.

Third and fourth, we evaluate cross-regime generalization on two real-data multi-output problems. The Yeast benchmark \cite{elisseeff2001yeast} is a 14-label multi-label classification task with 1,500 training examples and 917 validation examples. The RF1 benchmark \cite{spyromitros2016mtr} is an 8-target regression task with 4,108 training examples and 5,017 validation examples. In both cases, the tables report final task metrics aggregated over three seeds from runs scheduled for up to 120 epochs.

We also conduct four controlled NYUv2 studies on the fixed 50\% training subset used for the
ablation. The stop-gradient study compares canonical BPGS with a coupled variant for 60 epochs
over three seeds. The batch-size study uses batch sizes 4, 8, and 16, again with three seeds per
setting. The first-batch study varies only the loader seed (1001, 2001, 3001) while fixing the
model seed. The overhead study compares BPGS and Kendall for 60 epochs over three seeds, recording
per-epoch wall-clock time and peak CUDA memory.

Finally, we use three synthetic robustness studies. The pure loss-rescaling study applies post-hoc scale multipliers $\times 1$, $\times 10$, $\times 100$, and $\times 1000$ to test invariance to raw loss magnitude. The scale-stress study uses the same scale grid in a synthetic multi-task setting designed to probe degradation under scale imbalance. The heterogeneous mixed-stress study evaluates three regimes, denoted \emph{clean}, \emph{conflict}, and \emph{noisy}, to test robustness when tasks differ not only in scale but also in interaction structure.

To test whether normalization alone explains the rescaling result, we compare BPGS, Kendall, and
Kendall with L1-normalized precision weights on the pure-rescaling grid. This diagnostic uses seeds
42, 123, and 999 and is reported separately from the main rescaling table.

\subsection{Compared methods and metrics}
The exact comparison set depends on the benchmark. The main NYUv2 benchmark compares BPGS against Static, Kendall uncertainty weighting \cite{kendall2018}, UWSO, and Nash-MTL \cite{navon2022nash}. The NYUv2 ablation compares the four BPGS variants against Kendall. The Yeast and RF1 benchmarks compare BPGS against Static, Kendall, PCGrad \cite{yu2020pcgrad}, a GradNorm-inspired gradient-norm proxy \cite{chen2018gradnorm}, and UWSO. The pure loss-rescaling and scale-stress studies compare BPGS, Kendall, and UWSO, while the heterogeneous mixed-stress study compares BPGS, Kendall, PCGrad, and UWSO.

For NYUv2, we report segmentation mIoU, depth absolute relative error, depth RMSE, mean angular error for normals, normal accuracy within $11.25^\circ$, total validation loss, and in the main benchmark also $\Delta_M$ against the Static baseline \cite{liu2019mtan}. For Yeast, we report macro-F1, micro-F1, and Hamming accuracy. For RF1, we report $R^2$, RMSE, and MAE. For the synthetic stress studies, the primary statistic is macro score, with worst-task score used in the heterogeneous analysis and the stress diagnostics.

\subsection{Reporting protocol}
Unless stated otherwise, values are means $\pm$ standard deviations across seeds. The main NYUv2,
Yeast, RF1, rescaling, and heterogeneous-stress results use seeds 42, 43, and 44; the synthetic
scale-stress study in Table~5 uses 10 seeds. Synthetic stress results aggregate per-seed summaries,
whereas NYUv2, Yeast, and RF1 tables use the final recorded epoch. The anonymized supplementary
materials contain the code, configurations, and reproduction instructions.

Because checkpoint summaries are not uniformly available across methods, mixing best- and
final-checkpoint numbers would introduce method-dependent extraction rules. We therefore use final
epoch metrics for all paper tables. Final NYUv2 runs select checkpoints by \texttt{val/miou} in max
mode, while the ablation and controlled studies use the dataset default
\texttt{val/total\_loss} in min mode. The uncertainty update uses $g_\theta=100$; full settings are
in Appendix~\ref{app:reproducibility}.
